\begin{proposition}
    \label{prop:quitar_terminos}
    Sea $x=(x_n)_{n\geq1}$ una sucesión equidistribuida módulo $1$.  Si eliminamos
    un número finito de términos, la sucesión resultante sigue siendo equidistribuida módulo $1$.
\end{proposition}

\begin{proof}
    Sin pérdida de generalidad, supongamos que eliminamos los primeros $M$ términos de la sucesión 
    y sea $y=(y_n)_{n\geq1}$, definimos $y_n = x_{n+M}$ para cada $n \in \mathbb{N}$. 
    
    Sea $(a,b) \subset \mathbb{T}$,    
    \[
        A_N((a,b),y) = \card \left\{M+1 \leq n \leq N+M : \partfrac{x_n} \in (a,b)\right\}
    \]
    por lo que,
    \[
        A_N((a,b),y) = A_{N+M}((a,b),x) - A_M((a,b),x)
    \]

    Como $0 \leq A_M((a,b),x) \leq M$ y la sucesión $(x_n)$ es equidistribuida, 
    \[
        \lim_{N \to \infty} \frac{A_N((a,b),y)}{N} =
        \lim_{N \to \infty}  \left(\frac{N+M}{N} \, \frac{A_{N+M}((a,b),x)}{N+M} - \frac{A_M((a,b),x)}{N} \right) =
        b-a
    \]

    Por lo tanto, la sucesión $(y_n)$ es equidistribuida módulo $1$.
\end{proof}
